% Generated draft: requires mathematical and source review.
\documentclass[11pt]{article}
\usepackage[margin=25mm]{geometry}
\usepackage{fontspec}
\setmainfont{Latin Modern Roman}
\usepackage{amsmath,amssymb,mathtools,unicode-math}
\setmathfont{Latin Modern Math}
\usepackage{enumitem,needspace,xurl,hyperref}
\setlength{\parindent}{0pt}
\setlength{\parskip}{5pt}
\setlength{\emergencystretch}{3em}
\newcommand{\N}{\mathbb{N}}
\newcommand{\Z}{\mathbb{Z}}
\newcommand{\Q}{\mathbb{Q}}
\newcommand{\R}{\mathbb{R}}
\newcommand{\C}{\mathbb{C}}
\newcommand{\F}{\mathbb{F}}
\newcommand{\A}{\mathbb{A}}
\newcommand{\PP}{\mathbb{P}}
\newcommand{\E}{\mathbb{E}}
\newcommand{\eps}{\varepsilon}
\newcommand{\dd}{\,\mathrm{d}}
\begin{document}
\section*{UnsolvedMath ID: 11000023}
\textbf{Source problem code:} catalogue entry
\subsection{Definitions and mathematical statement}
Let \(\operatorname{Mod}_g\) be the mapping class group in the source's surface convention. Fix a finite generating set \(S\), and let \(\|x\|_S\) be the word length, the least number of letters from \(S\cup S^{-1}\) expressing \(x\). Under this interpretation of the unprovided norm convention, the problem is to prove
\[
 \exists C,K\geq0\quad\forall u,v\in\operatorname{Mod}_g:
 \quad u\text{ conjugate to }v
 \ \Longrightarrow\ 
 \exists h\in\operatorname{Mod}_g:\quad
 u=hvh^{-1},\quad
 \|h\|_S\leq K\max\{\|u\|_S,\|v\|_S\}+C.
\]
The constants depend only on the fixed data denoted \(S\) in the source and are independent of \(u,v\). The letter \(h\) replaces the printed conjugator \(g\) to avoid confusion with genus. No uniformity over all genera or all generating sets is asserted.

The source reports Masur--Minsky [MM2, Theorem 7.2] for pseudo-Anosov \(u,v\), and suggests extending hierarchy methods to the general case. It notes that this would still not answer a following basic problem, whose statement is absent. The excerpt does not define \(S\), the norm or the exact surface conventions; if \(S\) denotes the surface rather than a generating set, the metric and dependence must be clarified accordingly.
\subsection{Short English statement}
Find a conjugator whose word length is bounded linearly by the larger word length of two conjugate mapping classes, with constants fixed independently of those classes.
\subsection{Sources}
\begin{itemize}
\item \emph{Problems on Mapping Class Groups and Related Topics}, edited by Farb, Problem 3.2, ``Conjugator length bounds,'' supplied PDF page 31. \url{https://www.math.uchicago.edu/~farb/papers/mcgbook.pdf}. Recorded access 2026-07-29; extraction pdf-font-structure.
\item Masur--Minsky [MM2, Theorem 7.2], cited in the excerpt for the pseudo-Anosov case and hierarchy methods. The complete bibliography entry is absent.
\item Local UnsolvedMath registry, \texttt{data/unsolvedmath/problems.json}, permanent ID \(11000023\), exact source code \texttt{catalogue entry}; supplied OPEN-TRIAGE summary mentions conjugator-length work without a precise subsequent reference.
\end{itemize}
\end{document}
